\let\textcircled=\pgftextcircled
\chapter*{Conclusion g\'en\'erale}
\label{chap:conclusion}
\addstarredchapter{Conclusion g\'en\'erale}
\markboth{Conclusion g\'en\'erale}{}

\section*{Rappel de la probl\'ematique}
\initial{L}e travail partait d'une question : \`a quelles conditions un syst\`eme
peut-il ex\'ecuter une action corrective sans validation humaine pr\'ealable, sans que
le co\^ut d'une erreur d\'epasse le b\'en\'efice de la rapidit\'e ?

\section*{D\'emarche et r\'esultats}
La r\'eponse apport\'ee tient en une contrainte : \emph{la r\'eversibilit\'e est la
condition de l'autonomie}. [Reprendre les r\'esultats chiffr\'es du
chapitre~\ref{chap:implementation}.]

\section*{Apports}
\begin{itemize}
  \item \textbf{Apport conceptuel} --- le d\'eplacement de la question, d'une
        confiance dans la d\'etection vers une borne sur le co\^ut de l'erreur ;
  \item \textbf{Apport m\'ethodologique} --- la propagation d'une contrainte unique
        \`a tous les niveaux du mod\`ele, v\'erifiable dans chaque vue ;
  \item \textbf{Apport op\'erationnel} --- une plateforme fonctionnant hors ligne,
        dont chaque geste est opposable.
\end{itemize}

\section*{Limites}
\begin{enumerate}
  \item La mesure de r\'ef\'erence de la boucle de contr\^ole n'est pas persist\'ee : apr\`es
        un red\'emarrage, la boucle s'abstient sur les actions ant\'erieures.
  \item L'enrichissement du catalogue reste subordonn\'e \`a une validation humaine ;
        c'est le seul point de la cha\^ine o\`u une attente subsiste.
  \item Les actionneurs ont \'et\'e \'eprouv\'es en posture de simulation ; le
        raccordement \`a des \'equipements de production reste \`a mener.
  \item [\`A compl\'eter : limites du jeu de donn\'ees, du p\'erim\`etre, de la dur\'ee.]
\end{enumerate}

\section*{Perspectives}
\begin{enumerate}
  \item \textbf{Persistance des mesures de r\'ef\'erence}, pour que la boucle de
        contr\^ole survive \`a un red\'emarrage ;
  \item \textbf{Raccordement aux \'equipements du parc}, avec une phase d'observation
        en posture de simulation avant bascule ;
  \item \textbf{Apprentissage des seuils de d\'egradation} par cible, pour r\'eduire les
        annulations inutiles ;
  \item \textbf{Fed\'eration inter-CIRT}, pour partager les qualifications valid\'ees
        sans partager les donn\'ees d'incident.
\end{enumerate}
